\subsubsection{Discusión}
\label{sec:discusion-pruebas}

\paragraph{Hallazgos de la sesión escolar}
Los dibujos y la lectura de instrucciones recibieron una valoración favorable. La escritura en pantalla fue la actividad que menos alumnos consideraron fácil. Ambos docentes señalaron el tamaño de los textos y los íconos como un aspecto con margen de mejora, aunque la encuesta no identifica qué dificultó el trazo.

Los docentes valoraron que la calificación sugerida pueda revisarse antes de incorporarla al progreso. También señalaron la disponibilidad de internet y dispositivos como barreras para usar Graf-IA en su escuela. Entre los alumnos predominó la preferencia por combinar tableta y cuaderno.

\paragraph{Alcance de la evaluación}
Las 24 entregas de los ejercicios de la sesión permitieron una primera medición del OCR con escritura de primer grado (Sección~\ref{sec:resultados-cer-alumnos}). Comparando solo letras y espacios en la zona del ejercicio, el CER fue de 14.16\,\%: 12.95\,\% en el lienzo y 14.71\,\% en las fotografías del cuaderno. Con esta muestra se alcanza de forma preliminar la meta de 15\,\% del objetivo OE-03, aunque con poco margen.

Los errores se concentraron en letras de forma parecida, en las marcas del cuaderno unidas a las palabras y en el encuadre: 13 de las 15 fotografías incluyeron texto ajeno al ejercicio. Esto sugiere dos medidas:
\begin{itemize}
    \item guiar el encuadre de las fotografías o recortar la zona del ejercicio antes de enviarla al OCR;
    \item orientar al alumno sobre el espaciado entre letras.
\end{itemize}
Verificar la meta requiere más entregas, más ejercicios y una referencia revisada por más de una persona.

\paragraph{Limitaciones}
Los resultados describen a los 17 alumnos y dos docentes que participaron por conveniencia. La sesión careció de grupo de control y de mediciones antes y después. Las preguntas se leyeron en voz alta y la tendencia de los niños a elegir respuestas positivas en escalas de satisfacción pudo influir en las valoraciones \cite{read2006funtoolkit}.
